\section{Derivative Recovery on the Same Surface}
\label{sec:recovery-surface-fidelity-branch}

The derivative-recovery branch begins after a terminal object has been selected.
The positive-forward requirement is same-surface structure strong enough to
recover classical continuation.  The branch must prove a named continuation
criterion on the original Navier--Stokes surface and prove that its bound
yields a uniform relaunch interval crossing \(T_*\).  That is the recovery
payment; it cannot be supplied by assuming continuation-level regularity.

The recovery step has a precise burden. It has to remain on the same
Navier--Stokes surface. The recovered control must belong to the original
solution, use the original pressure-viscosity equation, and return a continuation
norm for the same branch. A recovery mechanism that succeeds on a different
carrier closes the original finite-time breakdown claim only after a same-surface
bridge is supplied.

The recovery target is criterion-explicit: a terminal packet, local profile,
scale decomposition, or geometric carrier must pay the full hypotheses of the
named continuation theorem actually invoked.  No unnamed regularity threshold
is allowed to stand in for that payment.

One recovery route uses packet control. If a finite cover or material atlas
carries uniform field control at positive scale, then finite-overlap patching can
convert local control into a global Sobolev readout. The obstruction is that the
packet has to be the same terminal packet produced by the original solution. A
packet that appears only after a limiting operation, a surface transfer, or a
change of carrier must be admitted back to the same-solution test.

Another recovery route uses geometric drain. If a carrier geometry produced a
coercive derivative identity, the required inequality has the form
\[
  \frac{d}{dt}\mathcal D(t)+c\mathcal H(t)\leq \mathrm{lower\ order\ terms},
\]
where \(\mathcal D\) controls derivatives and \(\mathcal H\) is a positive
higher-order dissipation. The stopping point is surface fidelity. A drain theorem
proved on an auxiliary geometric carrier has to be transported back to the flat
periodic or whole-space Navier--Stokes surface. The proof cannot spend geometric
curvature that the original equation does not provide.

A third recovery route uses endpoint or local-critical criteria. These criteria
can turn a local field readout into continuation, provided the readout is strong
enough and tied to the same terminal branch. The obstruction is again witness
fidelity: the endpoint readout has to be a readout of the selected object, not a
separate diagnostic profile.

The branch leaves a surface-fidelity residue. The residue is not a
single formula. It is a mismatch between the place where recovery is available
and the place where the breakdown claim lives. The manuscript tracks
this residue explicitly.

In CM-exit language, derivative recovery belongs mostly to the pass side. If the
terminal object keeps the required carrier, participation, and field readout,
then recovery supplies the continuation conclusion. If recovery fails because the
object lacks a positive carrier, loses the same equation, or has no positive-scale
field readout, the failure is classified through the corresponding membership
requirement.

This branch has two roles in the final paper. First, it explains the
classical continuation readout used by the pass branch. Second, it prevents
recovery mechanisms from being promoted before their surface-fidelity bridges
are supplied. The CM-exit method keeps that distinction sharp: recovery
belongs to the same-solution branch only after the selected terminal object has
passed the membership requirements needed for continuation.
